% GENERATED FILE. Edit canonical lessons, then run npm run generate:books.
% canonical-source-hash: fnv1a64:77f555ff6c56014d
% canonical-lessons: HI-C376-zahir, HI-C376-niscit, HI-C376-avasya, HI-C376-lazimi, HI-C376-ubharna

\chapter{Clearly So}
\label{ch:hi-clearly-so}
\hlchaptermodality{376}

\begin{chapteropening}
\textbf{By the end of this chapter:} \emph{I can say what is certain to happen if something is done.}

Complete the last lesson of chapter 376: I can say what is certain to happen if something is done.
\end{chapteropening}

\section[\texorpdfstring{\hi{ज़ाहिर} (zāhir)}{ज़ाहिर (zāhir)}]{\hi{ज़ाहिर} (zāhir) — evident, plain}

\label{lesson:HI-C376-zahir}



\begin{quote}
Before the new one: say the Hindi for a hint, then the Hindi for to charm.
\end{quote}

\subsection*{You'll want to know: \hi{ज़ाहिर}}
\textbf{\hi{ज़ाहिर}} — \emph{zāhir} — \textquotedblleft{}evident, plain\textquotedblright{}.

\begin{quote}
\hi{ज़ाहिर} \hi{है} \hi{कि} \hi{अगर} \hi{तुम} \hi{सोओगे} \hi{नहीं} \hi{तो} \hi{थक} \hi{जाओगे।}
\end{quote}

\begin{quote}
\emph{zāhir hai ki agar tum sooge nahīṃ to thak jāoge.}
\end{quote}

\begin{quote}
It's plain that if you don't sleep, you'll get tired.
\end{quote}

\begin{cousinweb}
\textbf{\hi{ज़ाहिर}} (\emph{zāhir}) is usually traced to Arabic \emph{ẓāhir}, \textquotedblleft{}outward, apparent\textquotedblright{}, which reached Hindi through Persian.
\end{cousinweb}

\begin{grammarlens}[title={What you've built}]
One more word for saying what happens if.
\end{grammarlens}

\subsection*{Your turn}
\begin{itemize}
\raggedright
  \item \emph{Say it:} \emph{zāhir}
  \item \emph{Say it:} \emph{zāhir}, once more
  \item \emph{Say it:} \emph{zāhir hai ki agar tum sooge nahīṃ to thak jāoge.} — It's plain that if you don't sleep, you'll get tired.
  \item \emph{Recall:} say the Hindi for a hint, then the Hindi for to charm, then say \emph{zāhir} again
\end{itemize}

\begin{tcolorbox}[breakable,colback=teal!4,colframe=teal!35!black,title={Before you move on}]
What does \hi{ज़ाहिर} mean? (\textquotedblleft{}Evident, plain\textquotedblright{}.) Say it once more.
\end{tcolorbox}
\section[\texorpdfstring{\hi{निश्चित} (niścit)}{निश्चित (niścit)}]{\hi{निश्चित} (niścit) — definite, fixed}

\label{lesson:HI-C376-niscit}



\begin{quote}
Before the new one: say the Hindi for to charm, then the Hindi for evident.
\end{quote}

\subsection*{You'll want to know: \hi{निश्चित}}
\textbf{\hi{निश्चित}} — \emph{niścit} — \textquotedblleft{}definite, fixed\textquotedblright{}.

\begin{quote}
\hi{अगर} \hi{तारीख़} \hi{निश्चित} \hi{होगी} \hi{तो} \hi{मैं} \hi{टिकट} \hi{ले} \hi{लूँगा।}
\end{quote}

\begin{quote}
\emph{agar tārīkh niścit hogī to maiṁ ṭikaṭ le lūṅgā.}
\end{quote}

\begin{quote}
If the date is fixed, I'll get a ticket.
\end{quote}

\begin{cousinweb}
\textbf{\hi{निश्चित}} (\emph{niścit}) is a tatsama, borrowed from Sanskrit \emph{niścita}, \textquotedblleft{}ascertained, settled\textquotedblright{} (Monier-Williams).
\end{cousinweb}

\begin{grammarlens}[title={What you've built}]
One more word for saying what happens if.
\end{grammarlens}

\subsection*{Your turn}
\begin{itemize}
\raggedright
  \item \emph{Say it:} \emph{niścit}
  \item \emph{Say it:} \emph{niścit}, once more
  \item \emph{Say it:} \emph{agar tārīkh niścit hogī to maiṁ ṭikaṭ le lūṅgā.} — If the date is fixed, I'll get a ticket.
  \item \emph{Recall:} say the Hindi for to charm, then the Hindi for evident, then say \emph{niścit} again
\end{itemize}

\begin{tcolorbox}[breakable,colback=teal!4,colframe=teal!35!black,title={Before you move on}]
What does \hi{निश्चित} mean? (\textquotedblleft{}Definite, fixed\textquotedblright{}.) Say it once more.
\end{tcolorbox}
\section[\texorpdfstring{\hi{अवश्य} (avaśya)}{अवश्य (avaśya)}]{\hi{अवश्य} (avaśya) — without fail, certainly (formal)}

\label{lesson:HI-C376-avasya}



\begin{quote}
Before the new one: say the Hindi for evident, then the Hindi for definite.
\end{quote}

\subsection*{You'll want to know: \hi{अवश्य}}
\textbf{\hi{अवश्य}} — \emph{avaśya} — \textquotedblleft{}without fail, certainly (formal)\textquotedblright{}.

\begin{quote}
\hi{अगर} \hi{समय} \hi{मिलेगा} \hi{तो} \hi{मैं} \hi{अवश्य} \hi{आऊँगा।}
\end{quote}

\begin{quote}
\emph{agar samay milegā to maiṁ avaśya āūṅgā.}
\end{quote}

\begin{quote}
If I get the time, I'll come without fail.
\end{quote}

\begin{cousinweb}
\textbf{\hi{अवश्य}} (\emph{avaśya}) is a tatsama, borrowed from Sanskrit \emph{avaśyam}, \textquotedblleft{}necessarily, inevitably\textquotedblright{} (Monier-Williams).
\end{cousinweb}

\begin{grammarlens}[title={What you've built}]
One more word for saying what happens if.
\end{grammarlens}

\subsection*{Your turn}
\begin{itemize}
\raggedright
  \item \emph{Say it:} \emph{avaśya}
  \item \emph{Say it:} \emph{avaśya}, once more
  \item \emph{Say it:} \emph{agar samay milegā to maiṁ avaśya āūṅgā.} — If I get the time, I'll come without fail.
  \item \emph{Recall:} say the Hindi for evident, then the Hindi for definite, then say \emph{avaśya} again
\end{itemize}

\begin{tcolorbox}[breakable,colback=teal!4,colframe=teal!35!black,title={Before you move on}]
What does \hi{अवश्य} mean? (\textquotedblleft{}Without fail, certainly (formal)\textquotedblright{}.) Say it once more.
\end{tcolorbox}
\section[\texorpdfstring{\hi{लाज़िमी} (lāzimī)}{लाज़िमी (lāzimī)}]{\hi{लाज़िमी} (lāzimī) — bound to happen, inevitable}

\label{lesson:HI-C376-lazimi}



\begin{quote}
Before the new one: say the Hindi for definite, then the Hindi for without fail.
\end{quote}

\subsection*{You'll want to know: \hi{लाज़िमी}}
\textbf{\hi{लाज़िमी}} — \emph{lāzimī} — \textquotedblleft{}bound to happen, inevitable\textquotedblright{}.

\begin{quote}
\hi{अगर} \hi{तुम} \hi{बहुत} \hi{ज़्यादा} \hi{काम} \hi{करोगे} \hi{तो} \hi{थकान} \hi{होना} \hi{लाज़िमी} \hi{है।}
\end{quote}

\begin{quote}
\emph{agar tum bahut zyādā kām karoge to thakān honā lāzimī hai.}
\end{quote}

\begin{quote}
If you work far too much, you're bound to get tired.
\end{quote}

\begin{cousinweb}
\textbf{\hi{लाज़िमी}} (\emph{lāzimī}) is usually traced to Persian \emph{lāzimī}, \textquotedblleft{}necessary\textquotedblright{}, which is built on Arabic \emph{lāzim}, \textquotedblleft{}binding, necessary\textquotedblright{}.
\end{cousinweb}

\begin{grammarlens}[title={What you've built}]
One more word for saying what happens if.
\end{grammarlens}

\subsection*{Your turn}
\begin{itemize}
\raggedright
  \item \emph{Say it:} \emph{lāzimī}
  \item \emph{Say it:} \emph{lāzimī}, once more
  \item \emph{Say it:} \emph{agar tum bahut zyādā kām karoge to thakān honā lāzimī hai.} — If you work far too much, you're bound to get tired.
  \item \emph{Recall:} say the Hindi for definite, then the Hindi for without fail, then say \emph{lāzimī} again
\end{itemize}

\begin{tcolorbox}[breakable,colback=teal!4,colframe=teal!35!black,title={Before you move on}]
What does \hi{लाज़िमी} mean? (\textquotedblleft{}Bound to happen, inevitable\textquotedblright{}.) Say it once more.
\end{tcolorbox}
\section[\texorpdfstring{\hi{उभरना} (ubharnā)}{उभरना (ubharnā)}]{\hi{उभरना} (ubharnā) — to emerge, to come to the surface}

\label{lesson:HI-C376-ubharna}



\begin{quote}
Before the new one: say the Hindi for without fail, then the Hindi for bound to happen.
\end{quote}

\subsection*{You'll want to know: \hi{उभरना}}
\textbf{\hi{उभरना}} — \emph{ubharnā} — \textquotedblleft{}to emerge, to come to the surface\textquotedblright{}.

\begin{quote}
\hi{अगर} \hi{पानी} \hi{कम} \hi{होगा} \hi{तो} \hi{नदी} \hi{के} \hi{पत्थर} \hi{उभरेंगे।}
\end{quote}

\begin{quote}
\emph{agar pānī kam hogā to nadī ke patthar ubhreṅge.}
\end{quote}

\begin{quote}
If the water goes down, the stones in the river will show.
\end{quote}

Past (it …): \hi{वह} \hi{सवाल} \hi{फिर} \hi{उभरा}, \emph{vah savāl phir ubhrā}.

\begin{cousinweb}
\textbf{\hi{उभरना}} (\emph{ubharnā}) is a tadbhava, inherited from Sanskrit \emph{udbharati}, \textquotedblleft{}rises up\textquotedblright{} (CDIAL 2038).
\end{cousinweb}

\begin{grammarlens}[title={What you've built}]
One more word for saying what happens if.
\end{grammarlens}

\subsection*{Your turn}
\begin{itemize}
\raggedright
  \item \emph{Say it:} \emph{ubharnā}
  \item \emph{Say it:} \emph{ubharnā}, once more
  \item \emph{Say it:} \emph{agar pānī kam hogā to nadī ke patthar ubhreṅge.} — If the water goes down, the stones in the river will show.
  \item \emph{Recall:} say the Hindi for without fail, then the Hindi for bound to happen, then say \emph{ubharnā} again
\end{itemize}

\begin{tcolorbox}[breakable,colback=teal!4,colframe=teal!35!black,title={Before you move on}]
What does \hi{उभरना} mean? (\textquotedblleft{}To emerge, to come to the surface\textquotedblright{}.) Say it once more.
\end{tcolorbox}
